% Part: applied-modal-logics
% Chapter: epistemic-logic
% Section: truth-at-w

\documentclass[../../../include/open-logic-section]{subfiles}

\begin{document}

\olfileid{aml}{el}{trw}

\olsection{الصدق عند عالم}

يعين كل نموذج معرفي، كما في المنطق الجهوي النظامي، ال!!{formula}s الصادقة
عند كل عالم منه. ونعبر عن أصل الدلالة العلائقية بقولنا: «يصدق
النموذج~$\mModel{M}$ !!{formula}~$!A$ عند العالم~$\OLId{latin-ordinary}{w}$». وتحد علاقة الصدق
استقرائيًا؛ وهي في المؤثرات غير الجهوية عين ما تقدم في المنطق الجهوي
النظامي.

\begin{defn}\ollabel{defn:mmodels}
  \emph{صدق !!a{formula}~$!A$ عند~$\OLId{latin-ordinary}{w}$} في~$\mModel M$، ورمزه
  $\mSat{\OLId{latin-looped}{M}}{!A}[\OLId{latin-ordinary}{w}]$، يحد استقرائيًا بما يأتي:
  \begin{enumerate}
  \tagitem{prvFalse}{\indcase{!A}{\lfalse}{لا يتحقق أبدًا
    $\mSat{\OLId{latin-looped}{M}}{\lfalse}[\OLId{latin-ordinary}{w}]$}.}{}
  \tagitem{prvTrue}{\indcase{!A}{\ltrue}{يتحقق دائمًا
    $\mSat{\OLId{latin-looped}{M}}{\ltrue}[\OLId{latin-ordinary}{w}]$}.}{}
  \item $\mSat{\OLId{latin-looped}{M}}{\OLId{latin-ordinary}{p}}[\OLId{latin-ordinary}{w}]$ إذا وفقط إذا كان~$\OLId{latin-ordinary}{w} \in \OLId{latin-looped}{V}(\OLId{latin-ordinary}{p})$.
  \tagitem{prvNot}{\indcase{!A}{\lnot !B}{$\mSat{\OLId{latin-looped}{M}}{\indfrm}[\OLId{latin-ordinary}{w}]$ إذا
    وفقط إذا كان~$\mSat/{\OLId{latin-looped}{M}}{!B}[\OLId{latin-ordinary}{w}]$}.}{}
  \tagitem{prvAnd}{\indcase{!A}{(!B \land !C)}{$\mSat{\OLId{latin-looped}{M}}{\indfrm}[\OLId{latin-ordinary}{w}]$
    إذا وفقط إذا كان~$\mSat{\OLId{latin-looped}{M}}{!B}[\OLId{latin-ordinary}{w}]$ و$\mSat{\OLId{latin-looped}{M}}{!C}[\OLId{latin-ordinary}{w}]$}.}{}
  \tagitem{prvOr}{\indcase{!A}{(!B \lor !C)}{$\mSat{\OLId{latin-looped}{M}}{\indfrm}[\OLId{latin-ordinary}{w}]$
    إذا وفقط إذا كان~$\mSat{\OLId{latin-looped}{M}}{!B}[\OLId{latin-ordinary}{w}]$ أو~$\mSat{\OLId{latin-looped}{M}}{!C}[\OLId{latin-ordinary}{w}]$}
    (أو كلاهما).}{}
  \tagitem{prvIf}{\indcase{!A}{(!B \lif !C)}{$\mSat{\OLId{latin-looped}{M}}{\indfrm}[\OLId{latin-ordinary}{w}]$
    إذا وفقط إذا كان~$\mSat/{\OLId{latin-looped}{M}}{!B}[\OLId{latin-ordinary}{w}]$ أو~$\mSat{\OLId{latin-looped}{M}}{!C}[\OLId{latin-ordinary}{w}]$}.}{}
  \tagitem{prvIff}{\indcase{!A}{(!B \liff !C)}{$\mSat{\OLId{latin-looped}{M}}{\indfrm}[\OLId{latin-ordinary}{w}]$
    إذا وفقط إذا تحقق كل من~$\mSat{\OLId{latin-looped}{M}}{!B}[\OLId{latin-ordinary}{w}]$ و$\mSat{\OLId{latin-looped}{M}}{!C}[\OLId{latin-ordinary}{w}]$، أو
    لم يتحقق أي من~$\mSat{\OLId{latin-looped}{M}}{!B}[\OLId{latin-ordinary}{w}]$ و$\mSat{\OLId{latin-looped}{M}}{!C}[\OLId{latin-ordinary}{w}]$}.}{}
  \item\ollabel{defn:sub:mmodels-box}
    \indcase{!A}{\Knows_a !B}{$\mSat{\OLId{latin-looped}{M}}{\indfrm}[\OLId{latin-ordinary}{w}]$ إذا وفقط إذا كان
    $\mSat{\OLId{latin-looped}{M}}{!B}[\OLId{latin-ordinary}{w}']$ لكل~$\OLId{latin-ordinary}{w}' \in \OLId{latin-looped}{W}$ بحيث~$\OLId{latin-looped}{R}_\OLId{latin-ordinary}{a} ww'$}
  \end{enumerate}
\end{defn}

وهنا ينبغي تقييد علاقات النفاذ. فإن مقتضى
\olref{defn:sub:mmodels-box} أن تصدق !!a{formula}~$\Knows_\OLId{latin-ordinary}{a} !B$ عند~$\OLId{latin-ordinary}{w}$
إذا لم يوجد~$\OLId{latin-ordinary}{w}'$ يحقق~$\OLId{latin-looped}{R}_\OLId{latin-ordinary}{a} ww'$. وهذا بعينه حكم المنطق الجهوي النظامي:
إذا لم يكن للعالم لاحق صدقت عنده كل صيغ~$\Box$ صدقًا خاليًا. غير أن حمل
$\Knows$ على \emph{المعرفة} يجعل هذا بعيدًا عن المراد؛ إذ لا نريد حالًا
يعلم فيها الفاعل~$!A$ و$\lnot !A$ معًا.

ومن وجوه دفع ذلك أن نجعل علاقة النفاذ في منطق المعرفة \emph{انعكاسية}
أبدًا؛ فيكون العالم الفعلي متسقًا مع معلومات الفاعل. ولهذا يستعمل في منطق
المعرفة عادةً~S5، وإن جاز العدول إلى أنساق أضعف بحسب المعنى الدقيق المراد
من علاقة~$\Knows_\OLId{latin-ordinary}{a}$.

\begin{figure}
  \begin{center}
    \begin{tikzpicture}[modal]
      \node[world] (w1) [label={[align=right]right:\mTrue{p}\\ \mFalse{q}}]{$\OLId{latin-ordinary}{w}_\OLMathNumeral{1}$};
      \node[world] (w2) [label={[align=right]right:\mTrue{p}\\ \mTrue{q}},
        above right=of w1]{$\OLId{latin-ordinary}{w}_\OLMathNumeral{2}$};
      \node[world] (w3) [label={[align=right]right:\mFalse{p}\\ \mFalse{q}},
        below right=of w1] {$\OLId{latin-ordinary}{w}_\OLMathNumeral{3}$};
      \draw[<->] (w1) to node {$\OLId{latin-ordinary}{a}$} (w2);
      \draw[->,loop left] (w1) to node {$\OLId{latin-ordinary}{a},\OLId{latin-ordinary}{b}$} (w1);
      \draw[<->] (w1) to [swap] node {$\OLId{latin-ordinary}{b}$} (w3);
      \draw[->,loop above] (w2) to node {$\OLId{latin-ordinary}{a},\OLId{latin-ordinary}{b}$} (w2) ;
      \draw[->,loop below] (w3) to node {$\OLId{latin-ordinary}{a},\OLId{latin-ordinary}{b}$} (w3) ;
    \end{tikzpicture}
  \end{center}
  \caption{نموذج معرفي بسيط.}
  \ollabel{fig:simple}
\end{figure}

\begin{prob}
  حدّد أي العبارات الآتية تصح في~\olref[aml][el][trw]{fig:simple}:
  \begin{enumerate}
    \item $\mSat{\OLId{latin-looped}{M}}{\lnot \OLId{latin-ordinary}{q}}[\OLId{latin-ordinary}{w}_\OLMathNumeral{1}]$؛
    \item $\mSat{\OLId{latin-looped}{M}}{\Knows_\OLId{latin-ordinary}{a} \lnot \OLId{latin-ordinary}{q}}[\OLId{latin-ordinary}{w}_\OLMathNumeral{1}]$؛
    \item $\mSat{\OLId{latin-looped}{M}}{\Knows_\OLId{latin-ordinary}{b} \lnot \OLId{latin-ordinary}{q}}[\OLId{latin-ordinary}{w}_\OLMathNumeral{1}]$؛
    \item $\mSat{\OLId{latin-looped}{M}}{\Knows_\OLId{latin-ordinary}{b} \OLId{latin-ordinary}{q} \lor \Knows_\OLId{latin-ordinary}{b} \lnot \OLId{latin-ordinary}{q}}[\OLId{latin-ordinary}{w}_\OLMathNumeral{2}]$؛
    \item $\mSat{\OLId{latin-looped}{M}}{\Knows_\OLId{latin-ordinary}{a}( \Knows_\OLId{latin-ordinary}{b} \OLId{latin-ordinary}{q} \lor \Knows_\OLId{latin-ordinary}{b} \lnot \OLId{latin-ordinary}{q})}[\OLId{latin-ordinary}{w}_\OLMathNumeral{2}]$؛
    \item $\mSat{\OLId{latin-looped}{M}}{\EKnows_{\{\OLId{latin-ordinary}{a},\OLId{latin-ordinary}{b}\}} \lnot \OLId{latin-ordinary}{q}}[\OLId{latin-ordinary}{w}_\OLMathNumeral{3}]$؛
    \end{enumerate}
\end{prob}

وإذا ثبت هذا الحد للصدق عند عالم، انتقلت إليه سائر المفاهيم الدلالية من
المنطق الجهوي النظامي، كالصلاحية الجهوية والاستلزام، مع حمل المؤثرات
الجهوية على تأويلها المعرفي الجديد.

ويبقى شرط الصدق لمؤثر المعرفة المشتركة~$\CKnows_\OLId{latin-looped}{G}$. فقد تقدم في
\olref[sfr][rel][ops]{sec} أن \emph{الإغلاق المتعدي}~$\OLId{latin-looped}{R}^+$ لعلاقة~$\OLId{latin-looped}{R}$
يحد بأنه:
\begin{align*}
  \OLId{latin-looped}{R}^+ &= \bigcup_{ \OLId{latin-ordinary}{n} \in \mathbb{N},\ \OLId{latin-ordinary}{n} \ge \OLMathNumeral{1}} \OLId{latin-looped}{R}^\OLId{latin-ordinary}{n},
\intertext{حيث}
  \OLId{latin-looped}{R}^\OLMathNumeral{1} & = \OLId{latin-looped}{R} \text{ و}\\
  \OLId{latin-looped}{R}^{\OLId{latin-ordinary}{n}+\OLMathNumeral{1}} & = \Setabs{\tuple{\OLId{latin-ordinary}{x},\OLId{latin-ordinary}{z}}}{\exists \OLId{latin-ordinary}{y} (\OLId{latin-looped}{R}^\OLId{latin-ordinary}{n} xy \land Ryz)}.
\end{align*}
% تصحيح اصطلاحي (0486-I1): رُقّمت قوى العلاقة من 1 على الاصطلاح المعتاد من غير تغيير الإغلاق المتعدي المقصود.
فإذا كانت~$\OLId{latin-looped}{G}$ مجموعة الفاعلين، عرّفنا
$\OLId{latin-looped}{R}_\OLId{latin-looped}{G}=(\bigcup_{\OLId{latin-ordinary}{b} \in \OLId{latin-looped}{G}}\OLId{latin-looped}{R}_\OLId{latin-ordinary}{b})^+$ بأنه الإغلاق المتعدي لاتحاد علاقات نفاذهم.

\begin{defn}
إذا كان~$\OLId{latin-looped}{G}' \subseteq \OLId{latin-looped}{G}$، جعلنا~$\mSat{\OLId{latin-looped}{M}}{\CKnows_{\OLId{latin-looped}{G}'} !A}[\OLId{latin-ordinary}{w}]$ إذا
وفقط إذا كان~$\mSat{\OLId{latin-looped}{M}}{!A}[\OLId{latin-ordinary}{w}']$ لكل~$\OLId{latin-ordinary}{w}'$ بحيث~$\OLId{latin-looped}{R}_{\OLId{latin-looped}{G}'}ww'$.
\end{defn}

\end{document}
